\section{Desenvolvimento e implementação}\label{chap:implementacao-final}

\subsection{Arquitetura e divisão de responsabilidades}

A plataforma é uma aplicação modular executada em um processo Python com interface web, complementada por um componente JavaScript no navegador e serviços externos. A separação lógica entre apresentação, orquestração e acesso a dados não corresponde a microsserviços próprios ou a um backend REST independente. O Streamlit gerencia o ciclo de renderização e o estado da sessão.

O ponto de entrada local é \texttt{src/app/main.py}; o ponto de entrada web é \texttt{streamlit\_app.py}. Este último ajusta o caminho de importação, força armazenamento por sessão e autenticação por conta de serviço, carrega as configurações disponíveis nos Secrets e encaminha a execução à mesma função principal. A conta de serviço autentica a aplicação perante o Earth Engine, não o visitante \citep{GOOGLE_EE_SERVICE_ACCOUNT,STREAMLIT_CLOUD_SECRETS}.

\begin{figure*}[htbp]
\centering
\resizebox{0.96\textwidth}{!}{% Diagrama baseado em src/app e nas integrações observadas no código.
\begin{tikzpicture}[
  node distance=7mm and 9mm,
  box/.style={draw, rounded corners, align=center, minimum height=11mm,
    text width=37mm, font=\small, fill=green!6},
  external/.style={box, fill=blue!5},
  link/.style={-{Stealth[length=2mm]}, thick}
]
\node[box] (user) {Usuário};
\node[box, right=of user] (ui) {Interface Streamlit\\mapas e controles};
\node[box, right=of ui] (pipeline) {Orquestração Python\\\texttt{pipeline.py}};
\node[external, right=of pipeline] (gee) {Google Earth Engine\\Sentinel-2};
\node[external, above=of pipeline] (climate) {NASA POWER\\dados climáticos};
\node[external, above=of ui] (map) {Nominatim e Esri\\busca e mapa-base};
\node[box, below=of ui] (location) {Serviço e repositório\\de localizações};
\node[box, below=of pipeline] (export) {Exportação\\PDF, XLSX, CSV ZIP};
\node[box, below=of location] (sqlite) {Modo local:\\SQLite};
\node[box, left=of sqlite] (session) {Modo web:\\memória da sessão};
\node[box, left=of location] (geojson) {Transferência GeoJSON\\arquivo do usuário};
\draw[link,<->] (user) -- (ui);
\draw[link,<->] (ui) -- (pipeline);
\draw[link,<->] (pipeline) -- (gee);
\draw[link,<->] (pipeline.north) -- (climate.south);
\draw[link,<->] (ui.south) -- (location.north);
\draw[link,<->] (location.south) -- (sqlite.north);
\draw[link,<->] (location.south west) -- (session.north east);
\draw[link,<->] (geojson.east) -- (location.west);
\draw[link] (pipeline.south) -- (export.north);
\draw[link,<->] (ui.north) -- (map.south);
\end{tikzpicture}
}
\caption{Arquitetura implementada, com processamento no servidor Python, interação no navegador e serviços externos. Fonte: elaborado pelo autor a partir dos módulos do sistema.}
\label{fig:arquitetura-atual}
\end{figure*}

\begin{table*}[htbp]
\centering\small
\caption{Responsabilidades dos principais componentes implementados.}
\label{tab:componentes-final}
\begin{tabular}{p{0.25\textwidth}p{0.66\textwidth}}
\hline
\textbf{Componente} & \textbf{Responsabilidade e fronteira} \\ \hline
\texttt{main.py}, estilos & Controles, estado, validação de requisições, apresentação e ações de exportação. \\
\texttt{maps.py} & Geocodificação, mapas de seleção, conversão de geometrias, métricas, tiles e pré-carregamento de grades. \\
\texttt{map\_interaction.py}, HTML & Ponte Streamlit e mapa Leaflet com consulta local de valores raster. \\
\texttt{pipeline.py} & Coordenação da coleção, produtos por índice e contexto meteorológico. \\
\texttt{earth\_engine.py}, \texttt{time\_series.py} & Filtros, máscara, índices e redução espacial por aquisição. \\
\texttt{anomalies.py}, \texttt{climate.py} & Heurística temporal, tendência, requisições climáticas, tabela, gráfico e resumo. \\
Serviço e repositórios de localizações & Validação e operações de armazenamento SQLite ou em memória de sessão. \\
\texttt{export\_context.py}, exportadores & Snapshot normalizado e serialização em PDF, XLSX e CSV ZIP. \\ \hline
\end{tabular}
\end{table*}

O desenvolvimento incremental ampliou o núcleo de mapas e séries com análise multiíndice, exportações e armazenamento intercambiável. A implementação avaliada parte da revisão \texttt{b7c325c}, acrescida das correções de contratos e testes descritas na avaliação técnica.

\subsection{Busca geográfica e mapa de seleção}

A busca usa o endpoint público Nominatim \texttt{/search}, com texto normalizado, formato \texttt{jsonv2}, limite de um resultado e tempo limite de 10 s. O retorno contém descrição, latitude, longitude e caixa delimitadora. O zoom é estimado pela extensão dessa caixa e limitado de 1 a 20. A primeira resposta não comprova a identificação de uma propriedade rural; sua função é aproximar a navegação da região desejada.

O mapa de seleção é construído com Folium, tiles Esri World Imagery e ferramenta de polígono do Leaflet Draw. São habilitadas edição e remoção de desenhos. A análise utiliza uma área ativa obtida do evento de desenho; a existência de múltiplos desenhos no componente não implementa comparação automática entre talhões.

A conversão para \texttt{ee.Geometry.Polygon} verifica tipo, fechamento, coordenadas finitas, limites geográficos, quantidade de vértices e área planar não nula. Nesse caminho também são rejeitadas interseções próprias estritas entre segmentos não adjacentes. Isso não equivale a um validador topológico completo de polígonos, relações entre anéis e todos os casos degenerados. Área, perímetro e centroide analíticos são obtidos pelo Earth Engine e convertidos para hectares, quilômetros e coordenadas.

\subsection{Estado da interface e coerência entre entradas e saídas}

O estado distingue área ativa, parâmetros de análise, resultados por índice, grade de valores, URLs de tiles, snapshot e arquivos preparados. A inicialização evita sobrescrever valores em cada reexecução. O processamento apresenta estados de espera, execução, sucesso, ausência de imagens e erro. Falhas de autenticação impedem a continuidade da interface principal; falhas de clima são tratadas como indisponibilidade complementar pelo pipeline.

A aceitação de desenho, troca de área e busca bem-sucedida invalidam resultados associados à seleção anterior. A busca sem resultado preserva a seleção existente. A restauração de uma área é concluída somente após validação e obtenção das métricas; não chama o geocodificador nem inicia análise automática. Chaves de revisão do mapa e controle de eventos repetidos evitam reaplicar um polígono excluído ou perder desenhos durante reexecuções.

A alteração de geometria, período ou conjunto solicitado de índices invalida resultados, metadados, grades, tiles, snapshot e arquivos preparados. A alternância visual entre índices já calculados preserva esses produtos e não inicia consulta remota. A exportação também confere o snapshot antes da geração, aplicando a mesma preocupação de coerência entre entradas e saídas em uma fronteira adicional.

\subsection{Gerenciamento de localizações}

O serviço \texttt{SavedLocationsService} recebe nome e GeoJSON, aplica validações e coordena repositórios intercambiáveis. O registro contém identificador UUID, nome exibido, chave normalizada de comparação, geometria, coordenadas de enquadramento e datas UTC de criação e atualização. O nome é limitado a 100 caracteres; espaços são normalizados e a comparação utiliza \texttt{casefold}, sem distinguir caixa.

As coordenadas \texttt{center\_lat} e \texttt{center\_lon} persistidas correspondem ao ponto médio da caixa delimitadora, não ao centroide geométrico calculado no Earth Engine. A recuperação utiliza esse tipo de enquadramento para navegação, mas recalcula as métricas analíticas da área. Essa distinção evita confundir um centro de visualização com a coordenada da consulta climática.

No modo SQLite, uma tabela \texttt{saved\_locations} mantém GeoJSON serializado em texto, nomes e metadados. O caminho é configurável por \texttt{SAVED\_LOCATIONS\_DB\_PATH}; o padrão utiliza o diretório de dados do usuário. A inicialização é idempotente, usa \texttt{PRAGMA user\_version=1} e rejeita esquema futuro incompatível. Consultas parametrizadas, unicidade do nome normalizado e transações protegem o contrato de gravação. O armazenamento é da instância local e não oferece isolamento autenticado por pessoa.

No modo web, o repositório recebe um dicionário da sessão Streamlit e retorna cópias independentes. Não utiliza banco externo nem promete recuperação após encerramento ou perda da sessão. O limite do conjunto é verificado para que as localizações possam ser transferidas em um arquivo aceito pelo importador.

\begin{figure*}[htbp]
\centering
\resizebox{0.95\textwidth}{!}{% Recorte das classes reais de gerenciamento de localizações.
\begin{tikzpicture}[
  class/.style={draw, align=left, inner sep=3mm, font=\small, fill=green!6},
  link/.style={-{Stealth[length=2mm]}, thick}
]
\node[class] (service) at (0,0) {\textbf{SavedLocationsService}\\
  + save\_location(name, geojson)\\
  + list\_saved\_locations()\\
  + get\_saved\_location(id)\\
  + delete\_saved\_location(id)\\
  + import\_locations(payload)\\
  + export\_locations()};
\node[class, right=16mm of service] (repo) {\textbf{SavedLocationsRepository}\\
  + initialize()\\
  + insert\_location(record)\\
  + insert\_locations(records)\\
  + list\_locations()\\
  + get\_location(id)\\
  + delete\_location(id)};
\node[draw, rounded corners, align=center, right=12mm of repo,
  font=\small, fill=blue!5] (db) {SQLite\\saved\_locations};
\node[class, below=12mm of repo] (session) {\textbf{SessionLocationsRepository}\\
  + initialize()\\
  + insert\_location(record)\\
  + insert\_locations(records)\\
  + list\_locations()\\
  + get\_location(id)\\
  + delete\_location(id)};
\node[draw, rounded corners, align=center, right=12mm of session,
  font=\small, fill=blue!5] (memory) {Memória\\da sessão};
\draw[link] (service) -- node[above, font=\scriptsize] {local} (repo);
\draw[link] (service.south) |- node[below, font=\scriptsize, pos=0.75] {web} (session.west);
\draw[link] (repo) -- (db);
\draw[link] (session) -- (memory);
\end{tikzpicture}
}
\caption{Serviço e alternativas de armazenamento de localizações. Fonte: elaborado pelo autor a partir das classes implementadas.}
\label{fig:classes-localizacao}
\end{figure*}

A transferência aceita uma \texttt{Feature} ou \texttt{FeatureCollection}, contendo geometrias Polygon e nome em \texttt{properties.name}, limitada a 100 áreas e $2\times1024^2$ bytes. A importação valida o lote e rejeita duplicatas sem sobrescrever registros. No SQLite, a inserção é transacional; na sessão, a atualização ocorre após as verificações. IDs e datas são gerados novamente, não copiados do arquivo.

Há operações de exclusão no serviço e nos repositórios, porém não há ação correspondente na interface atual nem fluxo de renomeação. Também não são persistidos histórico de análises, credenciais ou objetos Earth Engine. A validação local de geometrias salvas verifica estrutura e área assinada de cada anel, mas não repete todos os testes de interseção do caminho de análise; as camadas não devem ser descritas como validadores topológicos equivalentes.

\begin{figure*}[htbp]
\centering
\resizebox{0.95\textwidth}{!}{\begin{tikzpicture}[
 box/.style={draw,rounded corners,align=center,text width=39mm,
 minimum height=13mm,font=\small,fill=green!5},
 link/.style={-{Stealth[length=2mm]},thick}]
\node[box] (draw) at (0,0) {Desenho Polygon\\validação e métricas};
\node[box] (active) at (5.1,0) {Área ativa em sessão\\GeoJSON e ee.Geometry\\resultados antigos limpos};
\node[box] (save) at (10.2,0) {Diálogo de nome\\serviço valida registro\\ID e datas gerados};
\node[box] (storage) at (15.3,0) {Repositório\\SQLite local\\ou memória da sessão};
\node[box] (restore) at (5.1,-2.6) {Restaurar: validar registro\\recriar geometria\\recalcular métricas e vista};
\node[box] (list) at (10.2,-2.6) {Listar e selecionar\\localização existente};
\node[box,fill=blue!5] (file) at (15.3,-2.6) {GeoJSON do usuário\\importação em lote\\ou exportação};
\draw[link] (draw)--(active);
\draw[link] (active)--(save);
\draw[link] (save)--(storage);
\draw[link] (storage.south west)--(list.north east);
\draw[link] (list)--(restore);
\draw[link] (restore)--(active);
\draw[link,<->] (storage)--(file);
\node[font=\small,align=center] at (0,-2.6) {Restauração não executa\\geocodificação nem análise};
\end{tikzpicture}
}
\caption{Demarcação, gravação e restauração da área ativa, com alternativas de armazenamento. Fonte: elaborado pelo autor a partir da interface e do serviço de localizações.}
\label{fig:fluxo-areas}
\end{figure*}

\subsection{Orquestração e análise multiíndice}

Os caminhos individual e multiíndice compartilham consulta de coleção e métodos de cálculo. A versão multiíndice busca a coleção, aplica a máscara e obtém o clima compartilhado antes de produzir os resultados de cada índice. Os nomes repetidos são removidos preservando ordem. O pipeline não executa treinamento, classificação de culturas ou diagnóstico causal.

Os objetos de imagem permanecem referências a operações Earth Engine, enquanto estatísticas e séries são transferidas ao Python. A redução de cada aquisição é recuperada por chamadas individuais em uma iteração; essa estratégia simplifica a estrutura de retorno, mas pode aumentar latência e quantidade de requisições. Não há fila de tarefas, processamento distribuído próprio ou mecanismo de retomada persistente de uma análise interrompida.

\subsection{Pré-carregamento e consulta interativa}\label{sec:hover-final}

A exploração espacial utiliza dois produtos relacionados, mas distintos: tiles coloridos para renderização e uma grade numérica para consulta. Após o processamento, \texttt{preload\_index\_values} escolhe a zona UTM pelo centro informado e cria projeção norte ou sul em escala nominal de 10 m. Cada composição é selecionada, reprojetada, recortada pelo polígono e convertida em ponto flutuante. Bandas \texttt{grid\_x} e \texttt{grid\_y} permitem localizar a origem dos índices da grade.

A extração usa \texttt{sampleRectangle}, que retorna uma matriz por banda na projeção de entrada \citep{GOOGLE_EE_SAMPLE}. A região é a caixa delimitadora projetada da área; valores mascarados recebem $-9999$. O código estima a quantidade de células considerando também duas bandas de coordenadas e recusa estimativas superiores a dois milhões de valores. Uma verificação posterior confere a dimensão efetivamente recebida dos índices. Esse limite é do cliente da aplicação, não garantia de que o serviço remoto aceitará qualquer requisição abaixo dele.

As matrizes são convertidas para \texttt{float32} little-endian, comprimidas com zlib e codificadas em Base64. O componente recebe CRS, transformação afim, primeira linha/coluna, dimensões, sentinela e valores. No navegador, \texttt{DecompressionStream("deflate")} recupera os bytes e \texttt{Float32Array} permite acesso por posição. A compatibilidade desses recursos e das bibliotecas carregadas por CDN é requisito de execução, ainda não avaliado em uma matriz de navegadores.

Os eventos de movimento e toque verificam pertencimento ao polígono por teste de cruzamentos, incluindo exclusão de anéis internos. O componente projeta longitude/latitude em UTM com Proj4 e usa os termos de escala e origem da transformação. Para uma grade sem rotação, a localização é:
\begin{align}
 c&=\left\lfloor\frac{x-x_0}{s_x}\right\rfloor-c_0,\\
 r&=\left\lfloor\frac{y-y_0}{s_y}\right\rfloor-r_0,\qquad q=r\,n_c+c.
\end{align}
O índice linear $q$ acessa a matriz do indicador ativo. Termos de rotação/cisalhamento não são utilizados pelo JavaScript; o procedimento pressupõe a grade alinhada construída na etapa anterior. Fora do polígono, fora da matriz, com valor não finito ou com a sentinela, são exibidas mensagens específicas. Valores válidos aparecem com quatro casas decimais.

Não ocorre chamada ao Earth Engine a cada movimento do cursor. A interação consulta a grade já transferida; os tiles de fundo e da camada temática ainda dependem de carregamento de rede. O valor consultado representa a célula reamostrada da composição mediana do período, não uma leitura instantânea de campo, uma aquisição individual, a média temporal ou a média espacial de toda a área.

Uma medição reproduzível de empacotamento utilizou três matrizes controladas de $128\times128$ células, cada uma com 65.536 bytes em float32. Os tamanhos zlib foram 89 bytes para a matriz constante, 845 para o gradiente repetido e 60.562 para valores pseudoaleatórios com semente 42; em Base64, resultaram 120, 1.128 e 80.752 bytes. O contraste mostra que compressibilidade depende dos dados e que Base64 acrescenta volume. Não foram medidos tráfego total, latência de produção ou número observado de chamadas evitadas; arquiteturalmente, o lookup não solicita valores remotos por evento.

\begin{figure*}[htbp]
\centering
\resizebox{0.96\textwidth}{!}{\begin{tikzpicture}[
 box/.style={draw,rounded corners,align=center,text width=41mm,
 minimum height=15mm,font=\small,fill=green!5},
 link/.style={-{Stealth[length=2mm]},thick}]
\node[box] (median) at (0,0) {Composições medianas\\por índice};
\node[box] (project) at (5.4,0) {Projeção UTM a 10 m\\recorte e bandas\\de coordenadas};
\node[box] (sample) at (10.8,0) {\texttt{sampleRectangle}\\matrizes, sentinela\\e limites de tamanho};
\node[box] (pack) at (16.2,0) {float32 + zlib + Base64\\CRS, transformação\\e origem da grade};
\node[box,fill=blue!5] (decode) at (16.2,-2.8) {Navegador: descomprimir\\Float32Array\\reuso pela chave da análise};
\node[box,fill=blue!5] (mouse) at (10.8,-2.8) {Movimento ou toque\\ponto dentro do polígono\\projeção por Proj4};
\node[box,fill=blue!5] (lookup) at (5.4,-2.8) {Linha e coluna\\índice linear da célula\\teste de limites e NoData};
\node[box,fill=blue!5] (tooltip) at (0,-2.8) {Tooltip: valor\\da composição mediana\\ou mensagem de ausência};
\draw[link] (median)--(project);
\draw[link] (project)--(sample);
\draw[link] (sample)--(pack);
\draw[link] (pack)--node[right,font=\scriptsize]{transferência inicial}(decode);
\draw[link] (decode)--(mouse);
\draw[link] (mouse)--(lookup);
\draw[link] (lookup)--(tooltip);
\node[font=\small,align=center] at (8.1,-4.5) {Movimento e toque consultam memória local, sem nova requisição de valor ao Earth Engine.\\Tiles cartográficos e analíticos são carregados por um caminho de rede separado.};
\end{tikzpicture}
}
\caption{Transferência da grade e consulta local por posição no navegador. Fonte: elaborado pelo autor a partir de \texttt{maps.py} e \texttt{index\_hover/index.html}.}
\label{fig:fluxo-consulta-pixel}
\end{figure*}

\subsection{Cores, legendas e mapas estáticos}

As paletas são sequenciais: branco--verde no NDVI, branco--azul no NDWI e tons claros--alaranjados--marrons no NDMI. No caminho interativo principal, os extremos são obtidos dos valores válidos da grade. Intervalos constantes recebem pequena expansão e conjuntos sem valores válidos usam $[-1,1]$. Os tiles são recortados pela geometria e obtidos com \texttt{getMapId}; a opacidade no Leaflet é 0,7.

A ponte Python transmite ao argumento \texttt{range} os mesmos extremos utilizados nos tiles. O navegador calcula mínimo, ponto médio e máximo da legenda a partir desse contrato, inclusive em intervalos negativos. Expansão de intervalos constantes e fallback sem dados são definidos antes da renderização, evitando escalas independentes entre camada e legenda.

No PDF, o renderizador estático usa intervalo fixo $[-1,1]$, thumbnail Earth Engine em EPSG:3857 e mosaico de tiles Esri. O contorno, gradiente e atribuição são desenhados com Pillow. Portanto, o PDF não é uma captura do componente interativo e pode ter contraste visual diferente. Falhas de rede ou renderização produzem estado explícito de mapa indisponível.

\subsection{Exportação e interoperabilidade}\label{sec:exportacao-final}

O snapshot concluído associa resultados à geometria e ao período. Antes de gerar um arquivo, a interface exige correspondência com os parâmetros atuais; o construtor normaliza datas e copia estruturas mutáveis, mantendo referências opacas a imagens e objetos gráficos. Não se realiza nova análise espectral para cada formato.

O PDF é gerado em memória com ReportLab em A4 paisagem, com uma página por índice presente no contexto, na ordem fornecida. Com os três índices, resulta em três páginas; com apenas um, em uma. Cada página reúne período, área, data de geração, média espacial da composição, tendência quando disponível, contagem de imagens quando fornecida, resumo meteorológico, série, anomalias e mapa estático. O nome da localização salva e coordenadas explícitas não são impressos automaticamente como identificação de propriedade. A lista de anomalias é resumida, exibindo até três ocorrências e indicação de excedentes. O PDF não substitui a saída tabular completa.

O XLSX contém \texttt{Resumo}, uma aba por índice selecionado, \texttt{Clima}, \texttt{Anomalias} e \texttt{Geografia}; são sete abas quando os três índices estão presentes. Nas abas espectrais, as colunas são \texttt{date}, \texttt{value}, \texttt{trend}, \texttt{anomaly} e \texttt{observation\_id}. A tendência é global à série. A marcação de anomalias utiliza o identificador de aquisição; registros legados sem identificador usam conjuntamente data e valor, não apenas data. A aba climática contém radiação normalizada, e a geográfica inclui área, centro, CRS e GeoJSON.

A coluna \texttt{classification} existe na saída de anomalias, mas o detector padrão não produz categoria agronômica. Os caminhos individual e multiíndice fornecem contagem, tendência e contexto climático sob contratos harmonizados, preservando ausência explícita quando um serviço não está disponível.

O ZIP é construído com \texttt{zipfile} e compressão DEFLATE. Contém \texttt{resumo.csv}, um CSV por índice, \texttt{clima.csv}, \texttt{anomalias.csv} e \texttt{geografia.csv}. Os arquivos usam UTF-8, separação por vírgula, cabeçalhos estáveis, datas ISO, booleanos textuais e campos vazios para ausência. Não são exportados rasters analíticos GeoTIFF nem todas as amostras por pixel. A transferência de localizações em GeoJSON é uma função separada da exportação de resultados.

\begin{figure*}[htbp]
\centering
\resizebox{0.96\textwidth}{!}{\begin{tikzpicture}[
 box/.style={draw,rounded corners,align=center,text width=39mm,
 minimum height=14mm,font=\small,fill=green!5},
 link/.style={-{Stealth[length=2mm]},thick}]
\node[box] (snapshot) at (0,0) {Snapshot de análise\\concluída com sucesso};
\node[box] (match) at (5.2,0) {Confere área, período\\e índices atuais\\se diferentes: reanalisar};
\node[box] (context) at (10.4,0) {Contexto normalizado\\datas, geometria\\e cópias dos dados};
\node[box] (pdf) at (15.6,2) {PDF por índice\\série e resumo locais\\mapa estático remoto};
\node[box] (xlsx) at (15.6,0) {XLSX\\abas compartilhadas\\e por índice};
\node[box] (zip) at (15.6,-2) {CSV UTF-8\\conjuntos tabulares\\ZIP DEFLATE};
\node[box] (download) at (20.8,0) {Bytes em memória\\download por formato\\sem persistir análise};
\draw[link] (snapshot)--(match);
\draw[link] (match)--(context);
\draw[link] (context)--(pdf);
\draw[link] (context)--(xlsx);
\draw[link] (context)--(zip);
\draw[link] (pdf)--(download);
\draw[link] (xlsx)--(download);
\draw[link] (zip)--(download);
\end{tikzpicture}
}
\caption{Geração sob demanda de arquivos a partir do snapshot de análise. Fonte: elaborado pelo autor a partir dos módulos de exportação.}
\label{fig:fluxo-exportacao}
\end{figure*}

Os formatos estruturados favorecem reutilização em ferramentas externas, inspeção dos resultados e transparência. Sua contribuição à reprodutibilidade é parcial: o contexto não inclui obrigatoriamente todos os IDs de imagens, versões de catálogo, parâmetros de reprojeção, máscaras efetivas e respostas remotas. A exportação de tabelas não constitui um pacote autossuficiente para repetir integralmente a análise.
